% ==============================================================
% Consolidated report, data0826 round.
% Part 1: the five compounds.  Part 2: the five spreadsheets shared
% for 5:3 FTCA and 7:3 FTCA.  Every dataset under both models.
% All figures/tables machine-generated by run_all.R.
% Descriptive only: no interpretation or recommendation is given.
% Built by run_all.R (do not edit output/report.tex; edit ./report.tex)
% ==============================================================
\documentclass[11pt]{article}
\usepackage[margin=1.9cm]{geometry}
\usepackage{amsmath}
\usepackage{amssymb}
\usepackage{graphicx}
\usepackage{booktabs}
\setlength{\parindent}{0pt}
\setlength{\parskip}{5pt}

\begin{document}

\begin{center}
{\Large\bfseries Aqueous $\mathrm{pKa}$ of Five PFAS Compounds}\\[3pt]
{\large Analysis with the Adjusted pH column, under two solvent models}\\[6pt]
Yu Cheng \quad August 26, 2026
\end{center}

\vspace{2mm}
\section*{1. What this report contains}

Every fit in this report uses the \emph{Adjusted pH} column of the spreadsheets in the
Box folder. Nothing else in the analysis was changed relative to the previous round.
Each dataset is analyzed under two solvent models and by two estimation routes, so each
dataset appears four times: two models $\times$ two routes.

\textbf{Part 1} covers the five compounds (5:3 FTCA, 7:3 FTCA, 8:2 FTCA, PFBA, PFOA),
one spreadsheet each. \textbf{Part 2} covers the five spreadsheets shared for 5:3 FTCA
and 7:3 FTCA (two CF2 resonances for 5:3; newest data, old data and a second CF2
resonance for 7:3). The two parts have the same internal structure, and every page
states the spreadsheet it was computed from.

\textbf{Solvent constants} (fixed, not read from the spreadsheets): dielectric constants
$\varepsilon_k = 63.50$, $58.71$, $53.91$ and water activities
$a_{\mathrm{H_2O},k} = 0.913$, $0.908$, $0.901$ for 40\%, 50\%, 60\% (v/v) ACN; pure
water $\varepsilon_w = 78.33$, $a_w = 1$.

\section*{2. The two models}

Each titration is described by the Henderson--Hasselbalch relation between the
$^{19}$F chemical shift and pH,
\begin{equation}
y_{k,i} \;=\;
\frac{\delta_{A,k} + \delta_{HA,k}\,10^{\,\mathrm{pKa}_k - \mathrm{pH}_{k,i}}}
     {1 + 10^{\,\mathrm{pKa}_k - \mathrm{pH}_{k,i}}} \;+\; e_{k,i},
\qquad e_{k,i} \sim N(0,\sigma^2),
\label{eq:hh}
\end{equation}
where $k = 1,2,3$ indexes the mixture, $\delta_{A,k}$ and $\delta_{HA,k}$ are the two
limiting shifts, and $\mathrm{pKa}_k$ is the effective $\mathrm{pKa}$ in mixture $k$.
The two models differ only in how $\mathrm{pKa}_k$ depends on the solvent and in the
point at which the aqueous value is read off:
\begin{align}
\text{Model 1 (Yasuda--Shedlovsky):}\qquad
\mathrm{pKa}_k &= B + A x_k - \log_{10} a_{\mathrm{H_2O},k},
\qquad x_k = 1/\varepsilon_k,
\label{eq:ys}\\[2pt]
\text{Model 2 (ACN concentration):}\qquad
\mathrm{pKa}_k &= \beta_0 + \beta_1 C_k,
\qquad C_k = 0.4,\ 0.5,\ 0.6 .
\label{eq:conc}
\end{align}
The aqueous value is read off at $x_w = 1/\varepsilon_w = 1/78.33 = 0.012767$ under
Model 1 and at $C = 0$ under Model 2. Model 2 carries no water-activity term.

\section*{3. How the estimates and their standard errors are formed}

Both models are fitted by two routes, and in both routes the aqueous value is a linear
functional of the two solvent-model coefficients. Writing $\theta$ for that pair of
coefficients ($\theta = (B, A)'$ under \eqref{eq:ys}, $\theta = (\beta_0, \beta_1)'$
under \eqref{eq:conc}) and $\hat V$ for its estimated covariance matrix,
\begin{equation}
\mathrm{pKa}^{\mathrm{aq}} \;=\; g'\theta,
\qquad
\widehat{\operatorname{Var}}\bigl(g'\hat\theta\bigr) \;=\; g'\hat V g
\;=\; \hat V_{11} + 2 g_2 \hat V_{12} + g_2^2 \hat V_{22},
\label{eq:quad}
\end{equation}
with the read-off vector
\begin{equation*}
g = (1,\ x_w)' = (1,\ 0.012767)' \ \text{ under \eqref{eq:ys}},
\qquad
g = (1,\ 0)' \ \text{ under \eqref{eq:conc}} .
\end{equation*}
Every standard error quoted in this report is an instance of \eqref{eq:quad}; the tables
give $\hat\theta$, its standard errors and its correlation, so each aqueous standard
error can be recomputed from the entries shown. Under Model 2 the second component of
$g$ is zero, so \eqref{eq:quad} collapses to $\hat V_{11}$: the aqueous
$\mathrm{pKa}$ is the intercept itself and its standard error is that of the intercept.
Under Model 1 the read-off point is not zero, so the cross term $2 x_w \hat V_{12}$
enters; $\hat A$ and $\hat B$ are strongly negatively correlated (the tables give the
value), and dropping that term would overstate the standard error.

\textbf{Route A --- joint fit (all titration points at once).} All $\sum_k n_k$
observations of \eqref{eq:hh} are fitted in a single non-linear least squares problem in
which $\mathrm{pKa}_k$ is replaced by \eqref{eq:ys} or \eqref{eq:conc}. The eight
parameters are $\theta$ together with the three pairs $(\delta_{A,k},\delta_{HA,k})$.
Here $\hat V$ is the $2\times2$ block of the fitted parameter covariance belonging to
$\theta$, computed as $\hat\sigma^2 \mathbf M^{-1}$ with
$\mathbf M = \sum_k w_k\, d_k d_k'$, $d_k = (1, x_k)'$ or $(1, C_k)'$, and $w_k$ the
information mixture $k$ retains about $\theta$ once its two limiting shifts are
profiled out. In the figures this route is the left panel and the blue star.

\textbf{Route B --- two-step.} Each mixture is first fitted on its own by three-parameter
non-linear least squares, giving $\widetilde{\mathrm{pKa}}_k$ with variance $\nu_k$;
these are the green squares and their error bars. The three values are then regressed on
$x_k$ (or $C_k$) by ordinary least squares,
$\tilde\theta = (\mathbf Z'\mathbf Z)^{-1}\mathbf Z' y$, where $\mathbf Z$ is the
$3 \times 2$ matrix whose rows are the $d_k'$, and
\begin{equation}
\tilde V_c \;=\; (\mathbf Z'\mathbf Z)^{-1}\mathbf Z'\,
\bigl[\operatorname{diag}(\nu_1,\nu_2,\nu_3) + \hat\tau^2 \mathbf I_3\bigr]\,
\mathbf Z\,(\mathbf Z'\mathbf Z)^{-1},
\qquad
\hat\tau^2 = \max\Bigl\{0,\ \mathrm{SS}_{\mathrm{res}} - \sum_k (1-h_k)\nu_k\Bigr\},
\label{eq:vc}
\end{equation}
where $\mathrm{SS}_{\mathrm{res}}$ is the residual sum of squares of the three points
about the fitted line and $h_k$ are its leverages. The component $\hat\tau^2$ measures
scatter of the three points beyond what their own measurement variances $\nu_k$ explain;
it is zero whenever the three points are consistent with the line, and then
\eqref{eq:vc} reduces to pure propagation of the $\nu_k$. The aqueous standard error of
this route, written $\mathrm{SE}_c$, is \eqref{eq:quad} evaluated at $\tilde V_c$. In
the figures this route is the right panel and the purple star with the black error bar.

\textbf{Lack-of-fit statistic.} $Q = \mathrm{SS}_{\mathrm{res}} / \sum_k (1-h_k)\nu_k$
compares the observed scatter of the three points about the line with the scatter their
measurement variances alone would produce. Under exact adherence to \eqref{eq:ys} or
\eqref{eq:conc} it follows a $\chi^2_1$ distribution; $\hat\tau^2 > 0$ exactly when
$Q > 1$. The tables report $Q$ and its $p$-value.

\textbf{Figure conventions.} Green squares: the per-mixture estimates
$\widetilde{\mathrm{pKa}}_k$ with $\pm 1.96\,$SE bars; in the Model 1 panels the known
offset $\log_{10} a_{\mathrm{H_2O},k}$ is added to each of them, so that the plotted
quantity is linear in $x_k$. Left panel:
the joint-fit line and the aqueous star with its $\pm 1.96\,$SE interval. Right panel:
the ordinary-least-squares line through the three points and the aqueous star with
$\pm 1.96\,\mathrm{SE}_c$. All intervals use the factor $1.96$.

\clearpage
\section*{Part 1 --- The five compounds}

\textbf{Aqueous $\mathrm{pKa}$, all five compounds under both models.}
\input{tab_overview_part1.tex}

\textbf{Lack-of-fit diagnostics.}
\input{tab_diag_part1.tex}

\input{body_part1.tex}

\clearpage
\section*{Part 2 --- The spreadsheets shared for 5:3 FTCA and 7:3 FTCA}

The five spreadsheets below are analyzed exactly as in Part 1. Two of them --- 5:3 FTCA
(main CF2) and 7:3 FTCA (newest) --- are the spreadsheets already used in Part 1 and are
repeated here so that all datasets of a compound can be read side by side.

\textbf{Aqueous $\mathrm{pKa}$, all five spreadsheets under both models.}
\input{tab_overview_part2.tex}

\textbf{Lack-of-fit diagnostics.}
\input{tab_diag_part2.tex}

\input{body_part2.tex}

\end{document}
